\documentclass[sigconf,nonacm]{acmart}
\usepackage{booktabs}
\usepackage{amsmath}
\usepackage{xcolor}
\settopmatter{printacmref=false}
\renewcommand\footnotetextcopyrightpermission[1]{}
\pagestyle{plain}
\newtheorem{finding}{Finding}

\title{A Clocked-Temporal Case Study of Post-Hoc Leakage in Agent-Memory Supersession Prediction}

\author{Anonymous}
\affiliation{\institution{Anonymous Institution}\city{Anonymous}\country{Anonymous}}

\begin{document}

\begin{abstract}
Supervised evaluation of agent-memory systems can fail in two distinct ways: temporal non-stationarity and post-hoc feature leakage. We separate these threats on a production deployment (2,635 records, 386 superseded). After removing a post-hoc counter \texttt{n\_superseded\_by} that matches the label with 98.4\% accuracy, random 5-fold CV yields F1 0.692, whereas fixed chronological splits yield 0.421--0.559. The absolute drop is 0.133--0.271; a sliding-window analysis varies from 0.344 to 0.921, showing that part of this gap is deployment-phase shift rather than future-information leakage. We therefore distinguish write-time constants, snapshot covariates, and post-write signals under an explicit as-of clock. A feature-level mutual-information audit shows that the leaked counter captures 81.3\% of target entropy, the best single honest feature captures 32.4\%, and a joint prediction vector captures 54.1\%; these are dataset diagnostics, not universal bounds. Three LoCoMo content tasks remain stable under temporal splitting, bounding our claim to event-dependent prediction. A governance comparison on the temporal hold-out shows that a transparent never-read baseline reviews 16.8\% of records and catches 93.9\% of future supersessions, outperforming a black-box model at the audited operating point on the reported metrics. We distill a clocked-temporal checklist for future event-dependent memory audits; we present this as a case study, not a field-wide solution.
\end{abstract}




\ccsdesc[500]{Computing methodologies~Artificial intelligence}
\ccsdesc[300]{Information systems~Data management systems}
\keywords{agent memory, temporal validation, supersession, evaluation methodology, data leakage}

\maketitle

\section{Introduction}
Agent-memory systems accumulate facts across sessions. A natural evaluation asks which memories will become obsolete, but such event prediction exposes two distinct threats: ordinary non-stationarity and \emph{post-hoc feature leakage}, in which a field computed after the outcome enters the feature set. These threats require different remedies, yet they are rarely separated.

Where these systems are evaluated by supervised prediction, the common practice is random or temporally uncontrolled resampling rather than an evaluation clock aligned with deployment. We therefore separate three quantities: (i) the apparent random-CV score, (ii) the deployment-aligned temporal score after removing label-derived fields, and (iii) the additional inflation caused by a post-hoc counter that directly encodes the label. Our scope is deliberately modest: we demonstrate the separation on one production memory deployment and use LoCoMo as a bounded negative control, not as evidence that every memory task is invalid.

We demonstrate these problems on production data from a multi-agent memory deployment and validate through cross-dataset analysis. The case study contributes five auditable results rather than a new field standard:
\begin{enumerate}
\item We quantify a deployment-sensitive gap: with the post-hoc counter removed, F1 drops from 0.692 under random CV to 0.421--0.559 under fixed chronological splits; sliding-window results show substantial cut-off sensitivity.
\item We use \texttt{n\_superseded\_by} as a controlled leak, not as a discovered deployment bug: it is a count of future supersessions, so its match with the label (98.4\% accuracy) is expected by construction. Its value is to provide an upper reference for how far an evaluation can drift when post-outcome fields are not separated.
\item A non-equivalent LoCoMo control shows that three content-facing tasks are stable under ordinary temporal splitting, delimiting the claim to event-dependent prediction.
\item A feature-level mutual-information audit shows that the best honest feature captures 32.4\% of target entropy, the controlled post-hoc counter captures 81.3\%, and a joint prediction vector captures 54.1\%. These are descriptive corpus diagnostics, not universal bounds.
\item We present a descriptive single-rater grouping of 386 superseded facts; 57.5% fall into routine software-development categories, but the grouping has no inter-rater validation.
\end{enumerate}

For future audits, we suggest a clocked-temporal checklist rather than a field standard: as-of timestamps, multiple rolling cut-offs, and event-prevalence and leakage diagnostics accompany every headline score.

\section{Background}
\label{sec:background}

\subsection{Agent Memory Systems}
Agent memory systems store and retrieve facts across sessions. Mem0 provides production-ready memory with scalable long-term storage. Zep tracks superseded facts with bi-temporal knowledge graphs. MemGPT pages context across hierarchical memory tiers. Engram implements a bi-temporal data model with a supersession chain, achieving 83.6% on LongMemEval-S. MemStrata uses deterministic (S,R,O) supersession to eliminate stale-fact errors without LLM calls. CALM decouples lifecycle preservation from query-time admission using learned gating policies. These systems differ in architecture, but their supervised event-oriented evaluations typically use random or otherwise temporally uncontrolled resampling; published system papers emphasize architecture and retrieval quality, while the deployment-aligned clock is seldom stated. We do not claim that every implementation commits this error. Instead, we make the hazard reproducible and provide a protocol for detecting it.

\subsection{Evaluation Benchmarks}
LongMemEval evaluates chat assistants on long-term memory across six categories with 500 questions. LoCoMo evaluates over ten long conversations with 1,986 questions. Both benchmarks measure retrieval accuracy: given a question, does the system retrieve the correct facts and produce a correct answer? Recent work extends evaluation to temporal validity (MemStrata) and lifecycle governance (CALM), but none addresses the question of whether the evaluation methodology itself is valid.

\subsection{Supersession}
When an agent writes a new fact that contradicts an existing memory about the same entity, the existing fact is superseded: marked as replaced but not deleted, preserving the audit trail. Supersession prediction asks: given a fact at write time, will it eventually be superseded? This is an event-based prediction task: the outcome depends on the future behavior of other agents (whether they write contradicting facts), not solely on the content of the fact itself. We show that this distinction has important implications for evaluation methodology.

\subsection{Temporal Split in Other Domains}
Temporal data splitting is well-established in time series forecasting, credit scoring, medical diagnosis, and recommendation systems. In credit scoring, temporal split is mandatory because economic conditions change over time and models trained on past data may not generalize to future market conditions. In medical diagnosis, temporal split prevents data leakage from evolving treatment protocols and changing patient demographics. In recommendation systems, temporal split accounts for changing user preferences and item availability. Agent memory evaluation faces analogous challenges: memory facts encode temporally-bound information that changes as the environment evolves, and random splits allow the model to learn from future facts that would not be available at deployment time.

The borrowed lesson is not simply ``do a temporal split,' but that each domain specifies the prediction clock and test-window origin. Rolling-origin forecasting repeats the origin over many windows~\cite{tashman2000,bergmeir2012,cerqueira2020}; cross-validation work warns that random folds are optimistic under serial dependence~\cite{bergmeir2012}; concept-drift research evaluates learners under explicit change regimes~\cite{gama2014}; and leakage research requires any post-outcome feature to be excluded or as-of-timestamped~\cite{kaufman2012,kapoor2023}. Agent memory combines these hazards: ordered writes, changing event rates, and metadata that can encode outcomes. We instantiate the known principles for agent-memory event prediction and quantify the overstatement.
\section{Methodology}
\label{sec:methodology}

\subsection{Dataset}
The dataset is derived from 30 days of production data from a multi-agent memory deployment serving 15 AI client instances through 22 adapters. Following anonymization, the benchmark comprises 2,635 samples: 386 labelled positive (superseded, 14.6\%) and 2,249 labelled negative (active or retired, 85.4\%). Each record includes content text, source agent identity, memory type, confidence score, entity count, quality value, read count, and timestamps.

\subsection{Features}
We evaluate 11 write-time features in three groups:
\begin{itemize}
\item Metadata (6): confidence, quality value, read count, entity count, type indicators (fact, decision)
\item Temporal (3): hour of day, day of week, days since recording
\item Content (2): content length, scope indicator
\end{itemize}
Entity count is computed by a built-in NER pipeline that identifies file paths (pattern matching on path separators), API endpoints (URL patterns), code identifiers (CamelCase and snake\_case patterns), and configuration parameters (key-value patterns). The pipeline processes content at write time and stores the count in the n\_entities field. We note that entity extraction quality varies by content type: code-related facts are more likely to contain recognizable entities than abstract observations.

We explicitly exclude \texttt{n\_superseded\_by} and \texttt{n\_supersedes}, which are post-hoc label-derived fields (Section~\ref{sec:leakage}). Table~\ref{tab:asof} makes the as-of semantics explicit.

\begin{table}[h]
\centering
\caption{As-of semantics of candidate features.}
\label{tab:asof}
\begin{tabular}{lp{4.2cm}}
\toprule
Feature & Valid prediction time \\
\midrule
Content length; type indicators; confidence; quality value; entity count; scope & Write time \\
Hour, weekday & Write time \\
Read count; use count & After read events; excluded from strict write-time scoring \\
Days since recording & Snapshot covariate; frozen at a declared clock \\
\texttt{n\_superseded\_by}; \texttt{n\_supersedes} & Never: post-hoc label-derived field \\
\bottomrule
\end{tabular}
\end{table}

\subsection{Model}
We use L2-regularized logistic regression with class weighting (balanced) for 500 epochs, learning rate 0.01, and L2 regularization of 0.01. Features are standardized before training. The temporal group requires an explicit as-of distinction. \emph{Days since recording} is a snapshot covariate: it is not fixed at the write event and is frozen at the declared prediction clock. \emph{Use count} and read count are behavioral covariates observed after writing; they are valid only for predictions made after those events and are excluded from strict write-time scoring. We therefore report the temporal protocol both for the full honest feature set and as a diagnostic sensitivity analysis, and the information decomposition treats these groups as post-write signals rather than write-time constants.

\subsection{Evaluation Protocol}
We compare two protocols:
\begin{itemize}
\item Random CV: stratified 5-fold cross-validation with random assignment.
\item Temporal split: train on the first $p$ percent of data chronologically, test on the remainder. We evaluate $p = 30, 40,$ and $50$. These are fixed ex ante sampling points of the event stream, not selected after observing test performance.
\end{itemize}
For random CV, F1 is macro-averaged over folds. For each fixed temporal split, F1 is reported with a 95\% percentile confidence interval from 100 iid bootstrap resamples of the hold-out predictions. The interval quantifies prediction instability conditional on the cut-off; it is not a population inference. Because hold-out predictions are temporally ordered, it is a conditional descriptive interval, not inferential evidence; the sliding-window analysis expresses sensitivity across cut-offs. F1 is a bounded score, so we describe leakage as a multiplicative ratio only descriptively and emphasize absolute differences. The chronological cuts are also power-limited: the 30\% and 40\% cut-offs leave roughly 90 and 98 training positives, respectively, and the 50\% cut-off leaves 198 test positives. Training support, test support, and event phase vary together; we interpret the range as a sensitivity boundary rather than three independent replications.

\paragraph{Why three split points.}
The three cut-offs control an unavoidable bias--power trade-off. A cut at 30\% leaves 296 positives for testing, maximizing power, but its training window precedes the late-September distribution shift; a cut at 50\% leaves the model more training data but only 198 test positives. The 40\% cut lies between them. Consistency across all three intervals is therefore more important than any single estimate: a genuine methodological effect should survive different allocations of the event horizon, whereas a split-specific artifact should not. We do not tune the split point to the test outcome.

\paragraph{Sample-size and validation caveats.}
With 14.6\% positives, the hold-out sets contain 198--296 positive examples. This scale is adequate for detecting the large observed drop, but it does not support tight subgroup comparisons. We therefore reserve per-agent and per-type analyses for descriptive inference and quantify the sensitivity of the causal entity-grounding check separately (Section~\ref{sec:entity}). The same fixed split rule, feature set, threshold, and seed are applied to every protocol; model selection uses the training partition only. Random CV is retained as a reference because it isolates distribution shift and is the prevailing practice; temporal split is the deployment-aligned estimate.

\section{Results}
\label{sec:results}

\subsection{A Controlled Post-Hoc Field}
\label{sec:leakage}
The feature n\_superseded\_by counts how many subsequent writes have superseded a fact. This directly encodes the prediction label: if n\_superseded\_by is greater than zero, the fact has been superseded. On our dataset, this feature matches the label with 98.4\% accuracy.

Including it inflates F1 from 0.692 to 0.930. The mutual information of the leaked feature is 0.488 bits (81.3\% of the target entropy H(Y) = 0.601 bits), compared with a maximum of 0.195 bits (32.4\%) for any honest feature. The gap between the leaked and honest univariate MIs is a direct diagnostic of how much post-hoc information the feature injects; it is not a bound on arbitrary multivariate predictors.

\subsection{Temporal Validation}
\label{sec:temporal}

Table~\ref{tab:temporal} shows supersession prediction under random CV and temporal splits. For random CV, the test-positives column reports the average positives per held-out fold.

\begin{table}[h]
\centering
\caption{Supersession prediction under random CV and temporal splits.}
\label{tab:temporal}
\begin{tabular}{lccc}
\toprule
Split & Test positives & F1 & 95\% CI \\
\midrule
Random 5-fold & 77 avg & 0.692 & [0.660, 0.720] \\
Temporal 30\% & 296 & 0.421 & [0.401, 0.438] \\
Temporal 40\% & 288 & 0.559 & [0.534, 0.593] \\
Temporal 50\% & 198 & 0.450 & [0.428, 0.475] \\
Majority class & --- & 0.461 & --- \\
\bottomrule
\end{tabular}
\end{table}

After removing the post-hoc field, temporal splits reduce F1 from 0.692 to 0.421--0.559; the absolute drop is 0.133--0.271, and the descriptive random/temporal ratio is 1.2--1.6. We interpret this as a deployment-generalization gap, not proof of hidden future-information leakage. Each temporal interval lies below the random-CV interval. The 0.138 spread between the 30\% and 40\% estimates warns against replacing one random score with one lucky temporal score. The 0.692 random score is the primary honest five-fold run; the 0.650 in Table~\ref{tab:ranking} is a separately constructed full feature set and is not the headline estimate. The majority-class baseline (F1 = 0.461) is close to the low-power split points, and the 40\% temporal model is 0.098 above it. Thus some signal survives temporal validation, but the honest estimate is far below random CV and is itself cut-off dependent.

An exploratory sliding-window analysis explains why a single split should not be over-read. Training on the first 20--55\% of the chronology and testing on the following 10\%, with a 5\% gap and a minimum of five test positives, yields F1 from 0.344 to 0.921. The windows were scanned after the main experiments; they are exploratory, uncorrected for multiple comparisons, and used only to motivate cut-off sensitivity. They do not establish a performance estimate. Two high-performing windows (F1 = 0.841 and 0.921) bracket the late-September burst, while a low-power window with only nine positives yields 0.344. The variation therefore reflects both temporal non-stationarity and event-density changes. Random CV fixes both hazards by mixing periods; temporal evaluation makes them visible. This is precisely why reporting only the maximum random-CV score is unsafe for an event task whose prevalence changes by nearly an order of magnitude across days.

\paragraph{Leakage versus shift.}
We do not infer ``leakage'' from the temporal drop alone. The post-hoc counter is a verified leakage channel because it is computed from the future label and matches the label on 98.4\% of records; removing it changes random-CV F1 from 0.930 to 0.692. By contrast, the remaining random-versus-temporal gap is compatible with non-stationarity, reduced training support, and prevalence changes. Our claim is therefore two-layered: (i) the post-hoc field is a leakage reference and is excluded from honest models; (ii) the temporal gap is a deployment-generalization gap whose size is conditional on the event history.

We then used LoCoMo (5,882 turns) as a non-equivalent control, not as an experimental manipulation of task type. Across three content-facing tasks---evidence identification (random 0.638 vs. temporal 0.604), speaker prediction (0.373 vs. 0.380), and temporal reference resolution (0.455 vs. 0.439)---the differences are within a few points. These tasks differ in labels and temporal structure, so the comparison cannot isolate task type causally; it shows only that our event effect is not a generic property of any supervised split.

\subsection{Feature Set Ranking}
\label{sec:ranking}

Temporal leakage does not merely reduce performance uniformly. It changes the ranking of competing feature sets.

\begin{table}[h]
\centering
\caption{Feature set ranking under random CV vs. temporal split.}
\label{tab:ranking}
\begin{tabular}{lccc}
\toprule
Feature set & Random F1 & Temporal F1 & Rank change \\
\midrule
With leaked features & 0.904 & 0.360 & down 1 to 2 \\
Metadata only & 0.557 & 0.483 & up 3 to 1 \\
Full (no leak) & 0.650 & 0.253 & down 2 to 3 \\
Content only & 0.509 & 0.242 & --- \\
Content extended & 0.508 & 0.241 & --- \\
\bottomrule
\end{tabular}
\end{table}

Under random CV, the leaked feature set ranks first; the 0.904 in Table~\ref{tab:ranking} is the mean over the ranking experiment, whereas the 0.930 in Section~\ref{sec:leakage} is the original five-fold run. Under temporal split, metadata only becomes the best performer. Features that encode post-write information improve random CV performance because they capture temporal patterns present in both training and test data. Under temporal split, these patterns differ between the training period and the test period, causing the complex features to generalize poorly. The ranking result does not show that temporal ranking is always ``correct'; it shows that a different ranking can arise once deployment time is respected, so a comparison made under random CV cannot be assumed stable.

\subsection{Mutual Information Analysis}
\label{sec:mi}

We estimate mutual information with an adaptive-binned discrete estimator over base-2 entropy on the full corpus. The chain-rule decomposition uses binary indicators for type, entity presence, and read presence; continuous fields are discretized by fixed bins. We apply no bias correction or resampling uncertainty, so all MI values are descriptive and should not be compared across datasets. The released script reports the exact bin construction. We compute the mutual information between the model predictions and the true outcome. The prediction vector achieves a joint empirical MI of 0.325 bits (54.1\% of H(Y) = 0.601 bits), a representation-level diagnostic rather than a bound on arbitrary predictors.

Per-feature mutual information analysis reveals the information content of each feature. The strongest honest feature is days since recording (I = 0.195 bits), followed by entity count (I = 0.135 bits) and use count (I = 0.112 bits); the leaked feature achieves 0.488 bits (81.3\%). Most features capture less than 15\% of the target entropy individually.

This is an empirical audit of the chosen representation, not a proof that every conceivable predictor is capped. The remaining 45.9\\% of entropy is therefore unexplained under the audited feature families, rather than irreducible in a universal mathematical sense.
Table~\ref{tab:feature-mi} separates the two failure modes that this aggregate number hides.

\begin{table}[h]
\centering
\caption{Univariate mutual information with the supersession outcome.}
\label{tab:feature-mi}
\begin{tabular}{lrr}
\toprule
Feature & MI (bits) & \% of $H(Y)$ \\
\midrule
\texttt{n\_superseded\_by} (leaked) & 0.488 & 81.3 \\
Days since recording & 0.195 & 32.4 \\
Entity count & 0.135 & 22.5 \\
Use count & 0.112 & 18.6 \\
Weekday & 0.090 & 14.9 \\
Quality value & 0.046 & 7.6 \\
Hour of day & 0.042 & 7.0 \\
Content length & 0.011 & 1.7 \\
Fact indicator & 0.011 & 1.8 \\
Confidence & 0.006 & 1.0 \\
Decision indicator & 0.006 & 1.0 \\
Shared scope & 0.002 & 0.4 \\
\bottomrule
\end{tabular}
\end{table}

The table warns against two different traps. First, \emph{days since recording} dominates the honest features, but it is a snapshot covariate rather than a strict write-time feature: its value changes as the evaluation clock moves toward the recorded event. It is legitimate for diagnostic analysis, yet a deployed prediction must declare the clock at which the feature is frozen. Second, confidence and quality carry almost no marginal information. Their apparent value in random models comes from correlations with post-write behavior, not from the stored score itself. A benchmark that reports only aggregate feature-group importance would leave both traps undetected.

\subsection{Information Decomposition}
\label{sec:info-decomp}

We decompose the mutual information to identify which feature groups contribute most to supersession prediction. Using the chain rule of mutual information:

I(Y; type) = 0.118 bits (19.7% of H(Y) = 0.601 bits)

I(Y; entities | type) = 0.131 bits (21.8% of H(Y))

I(Y; reads | type, entities) = 0.165 bits (27.4% of H(Y))

Total explained: 0.414 bits (68.9% of H(Y))

Residual: 0.187 bits (31.1% of H(Y))

The residual entropy (31.1%) represents information about supersession that cannot be captured by source, type, entity grounding, or read patterns. This information is contained in future agent behavior (which agent writes what, when) that is fundamentally unknowable at write time. For the audited feature families and the empirical chain-rule estimator, a Bayes-optimal predictor over exactly these variables cannot reduce uncertainty below 31\\% of the initial entropy. This is a dataset-specific decomposition, not a universal impossibility theorem.

Temporal leakage occurs because random CV artificially reduces this residual by allowing the model to learn from future patterns that will not be available in production.

\subsection{Entity Grounding and Memory Stability}
\label{sec:entity}

Entity-anchored facts are superseded at near-zero rates across all memory types, while entity-free facts are superseded at elevated rates.

\begin{table}[h]
\centering
\caption{Supersession rate by memory type and entity presence.}
\begin{tabular}{lrrrr}
\toprule
Type & No Ent \% & With Ent \% & Chi-sq & p \\
\midrule
Procedure & 99.2 & 0.0 & 116.3 & <0.001 \\
Experience & 58.2 & 0.0 & 227.7 & <0.001 \\
Incident & 26.9 & 0.0 & 36.5 & <0.001 \\
Fact & 17.8 & 0.2 & 108.1 & <0.001 \\
Decision & 11.3 & 0.0 & 16.4 & <0.001 \\
\bottomrule
\end{tabular}
\end{table}

The within-type tests are significant, but the large chi-square statistics partly reflect the sample sizes and near-zero cell counts for grounded records; we interpret the direction and stratified consistency, not the p-values alone. The odds ratio is 0.002: entity-anchored facts are much less likely to be superseded. This effect is not driven by a confound with memory type, as it holds within every type. However, cross-system validation using MemStrata rules shows that this effect does not generalize beyond entity-keyed supersession mechanisms. Within MemTether, the mechanism is structural: entity-keyed records can be deterministically matched for supersession, while entity-free records lack this coordination point. The coordination point is analogous to a primary key in databases: keyed records have unambiguous identity and can be updated deterministically; non-keyed records allow duplicates. When a new fact targets the same entity, the system detects the match and applies supersession. Entity-free facts lack this key, so new facts do not explicitly target them and semantic overlap goes undetected.

\subsection{Per-Agent Analysis}
\label{sec:per-agent}

Codex writes 1.6x more volatile facts than WorkBuddy (17.0\% vs. 10.9\%). The mechanism is entity grounding: Codex facts average 3.6 entities per fact (24.6\% with entities), while WorkBuddy facts average 14.0 entities (74.0\% with entities). Both agents produce facts with similar confidence scores and quality values. We have no independent gold standard for the NER pipeline, so lower coverage is an operational observation rather than a validated extraction error rate.

Improving entity extraction for Codex-style agents would reduce the supersession rate without restricting writes. Specifically, the NER pipeline should be extended to recognize technical entities (file paths, API endpoints, configuration parameters) that are currently missed in Codex-generated content. This intervention is directly actionable: it requires modifying the NER pipeline, not restricting agent behavior. The expected effect is a reduction in the supersession rate from 17.0% to approximately 11% (matching WorkBuddy), based on the assumption that entity grounding is the primary driver of the difference.

The per-type evidence identifies where the extraction repair should start. Procedures written by Codex have a 7.8\% entity-grounding rate but a 92.2\% supersession rate, whereas the ten procedures that do acquire an entity reference have no observed supersession. By contrast, configuration-related records are grounded 71.5\% of the time and are superseded only 12.4\% of the time. Thus the instability is not a generic Codex artifact. It concentrates in procedure-like records whose natural-language text names commands or file locations without emitting a machine-comparable key. Pattern-based recognition of shell commands, relative and absolute paths, service names, and registry or configuration keys would target exactly this residue. Because this is a stratified observational result, the estimate is directional rather than causal; a before--after deployment with unchanged write volume is the appropriate validation.


\subsection{Supersession Taxonomy}
\label{sec:taxonomy}

\begin{table}[h]
\centering
\caption{Descriptive single-rater grouping of 386 superseded facts.}
\label{tab:taxonomy}
\begin{tabular}{lrr}
\toprule
Grouping & Count & \% \\
\midrule
Iterative refinement & 113 & 29.3 \\
Tool usage notes & 73 & 18.9 \\
Other & 68 & 17.6 \\
Environment snapshots & 39 & 10.1 \\
Deployment records & 36 & 9.3 \\
Configuration decisions & 33 & 8.5 \\
Incident reports & 24 & 6.2 \\
\bottomrule
\end{tabular}
\end{table}

The first author manually assigned every superseded fact to one of seven categories after reading the full content of a purposive 100-fact sample and the category label of all 386 positives; no second independent coder was available. The labels are reproducible from the released annotation file, but they have no inter-rater agreement estimate and no adjudication protocol. We therefore use the taxonomy to explain event composition, not as a validated coding instrument. Iterative refinement (29.3\%) is the dominant cause: the agent iterates on a problem and each solution replaces the previous attempt. Tool usage notes (18.9\%) become stale when tools are updated. Environment snapshots (10.1\%) capture external service states. Iterative refinement, tool usage notes, and deployment records account for 57.5\% of supersessions. We interpret this as evidence that many supersession events are routine development activity rather than corruption, but the taxonomy is descriptive and should not be treated as a validated psychological construct.

\subsection{Survival Analysis}
\label{sec:survival}

The median time from write to supersession is 2.5 days. Supersession follows a burst pattern: 69.5\% of facts survive their first 24 hours, but only 4.5\% survive 72 hours. Facts are superseded fastest (median 0.3 days), while procedures persist longest (median 2.5 days).


\begin{table}[h]
\centering
\caption{Survival analysis: percentage of facts still active after N hours.}
\begin{tabular}{lr}
\toprule
Time & \% Still Active \\
\midrule
1 hour & 81.2 \\
6 hours & 76.4 \\
12 hours & 74.8 \\
1 day & 69.5 \\
2 days & 54.4 \\
3 days & 4.5 \\
7 days & 2.7 \\
\bottomrule
\end{tabular}
\end{table}
\subsection{Temporal Distribution of Supersession}
\label{sec:time-dist}
Supersession is not uniformly distributed across the deployment period.

\begin{table}[h]
\centering
\caption{Daily supersession rate.}
\begin{tabular}{lrrr}
\toprule
Date & Total & Superseded & Rate \\
\midrule
Sep 13 & 223 & 48 & 21.5 \\
Sep 14--23 & 858 & 49 & 5.7 \\
Sep 24 & 197 & 38 & 19.3 \\
Sep 25 & 189 & 164 & 86.8 \\
Sep 26 & 47 & 40 & 85.1 \\
Sep 27 & 111 & 34 & 30.6 \\
Sep 28 & 48 & 1 & 2.1 \\
Oct 01--17 & 745 & 1 & 0.1 \\
\bottomrule
\end{tabular}
\end{table}

The burst on September 25--27 corresponds to a period of intense development activity: 92 procedures were written by Codex in a single day (56.1% of that day’s supersessions). All 164 superseded facts on September 25 have zero entity references. After September 28, the supersession rate drops to 0%, corresponding to a shift from active development to testing and documentation.

This history has two methodological consequences. First, event prevalence is endogenous to operator behavior. A deployment that oscillates between configuration and documentation can change the positive rate by an order of magnitude within days, so a single random draw from all dates conflates task difficulty with workflow phase. Second, the post-burst silence is itself informative. Once routine configuration writes stopped, there were few competing statements to supersede; a low raw supersession rate after this date does not imply that older facts became intrinsically safer. Temporal evaluation therefore needs to report event prevalence, training-window phase, and censoring alongside F1. Without those covariates, a seemingly strong score may simply describe a quiet period.

\section{Discussion}
\label{sec:discussion}

This section separates what the case study observed from what we think it suggests. The observed items are the leakage reference, the random/temporal gap, the ranking change, the feature audit, the stratified entity association, the descriptive taxonomy, and the governance comparison at one cut-off. The suggestions are the checklist and monitoring triage. We label the latter explicitly and do not present them as validated field rules.

\subsection{Implications for Memory Governance}
The MI audit shows that the selected feature families leave much uncertainty unexplained. Governance should use post-write signals with an explicit as-of clock and a transparent baseline. Monitoring entity-grounding coverage is plausible, but agent identity conflates workflow, volume, and content mix; a defensible rule is to monitor low-coverage procedure-like records and compare any agent-differentiated threshold against the never-read baseline.


\subsection{Implications for Evaluation Practice}
For future event-dependent audits, we suggest a \emph{clocked-temporal checklist}, not a field standard. It asks for four disclosures: (1) an as-of timestamp for every feature; (2) rolling-origin or sliding-window validation over multiple cut-offs; (3) event prevalence, censoring, and training-window phase for every score; and (4) a leakage audit that reports post-hoc counters separately. Random CV may be reported only as a comparison, with one deployment-aligned score designated primary. This checklist makes mismatches measurable; it does not solve non-stationarity.



\subsection{Governance Strategy Comparison}
\label{sec:governance}

We compare seven governance strategies on the temporal hold-out set (1,318 facts, 198 superseded):

\begin{table}[h]
\centering
\caption{Governance strategy comparison on temporal hold-out set.}
\label{tab:governance}
\begin{tabular}{lrrrr}
\toprule
Strategy & Reviewed & Caught & Precision & Transparency \\
\midrule
No governance & 0\% & 0\% & --- & --- \\
Review all facts & 100\% & 100\% & 0.150 & High \\
Never-read heuristic & \textbf{16.8\%} & \textbf{93.9\%} & \textbf{0.842} & \textbf{High} \\
ML model & 77.5\% & 91.9\% & 0.178 & Low \\
Combined & 79.2\% & 98.5\% & 0.187 & Medium \\
\bottomrule
\end{tabular}
\end{table}

A simple monitoring heuristic---flagging facts that have never been read by the evaluation cut-off---is not a write-time predictor. It audits a post-write state at a declared review time, and its clock is the temporal test boundary; read counts after that boundary are not used. Because the base rate is 15.0\%, the comparison depends on review cost: it is a cost-sensitive comparison at selected points, not a formal Pareto test. The heuristic reviews 16.8\% of facts, catches 93.9\% of supersessions, and has precision 0.842; the model reviews 77.5\%, catches 91.9\%, and has precision 0.178. We report no full precision--recall or cost curve, so the result establishes a strong governance baseline rather than universal superiority. Because never-read status can correlate with low-attention clients or content types, deployment should monitor access bias and allow an operator override. The 6.1\% of missed supersessions are borne by later readers and by the operator who must adjudicate them; adversarial agents can also game the signal by reading their own low-quality writes. We therefore treat the heuristic as a triage tool for human review, not as an automatic deletion policy.

In this deployment, a transparent post-write signal is useful for governance triage, while a black-box prediction would add accountability costs without improving the audited catch rate. The never-read check is auditable in the narrow sense that an operator can verify the read counter for a flagged fact; this does not make the governance process fully auditable, because access telemetry, threshold choice, and overrides remain policy decisions.

Table~\ref{tab:governance} shows why the operating point is normative. Flagging procedures reviews 4.9\% of records with precision 0.953 but catches only 30.8\%; flagging entity-free records reaches 99.5\% recall at 83.5\% effort with base-rate precision. The never-read heuristic is not Pareto-optimal in every dimension, but it aligns review effort with access evidence and can be logged and explained.

The governance result should not be mistaken for a claim that machine learning is useless. A threshold-tuned model can be useful when recall requirements exceed what the heuristic provides or when read telemetry is incomplete. Rather, the comparison shows that a memory benchmark should include transparent access-based baselines. If a learned governance policy does not beat such a baseline while providing an audit story, the additional complexity is not justified by reported accuracy alone.



\subsection{Fairness, Accountability, and Transparency Implications}
\label{sec:faacct-implications}

Our findings support a concrete accountability model for memory governance, but we state its limits: there is no protected-attribute analysis and no user-facing harm study here. The contribution is a measurement discipline that prevents a deployment artifact from becoming the basis for system selection.

\textbf{Operational fairness.} We do not treat software agents as protected groups. The relevant fairness concern is benchmark selection: when a volatile phase or a low-quality post-hoc feature dominates a random split, a system can be adopted or rejected because of a measurement artifact. Our per-agent analysis illustrates the stakes rather than asserting disparate impact: Codex produces 1.6x more superseded facts than WorkBuddy, and the observed difference is accompanied by lower entity grounding (24.6\% versus 74.0\%). A fair evaluation must separate event prevalence, entity-keying implementation, and prediction skill before comparing systems.

\textbf{Accountability.} The never-read check is more auditable than an ML prediction because its input is a directly inspectable counter, whereas model decisions require trust in learned representations. This is a relative property, not a claim of complete auditability: telemetry quality, threshold policy, and operator overrides must also be recorded. Our comparison supports transparent triage at the audited thresholds; mapping to a particular legal regime remains domain-specific.

\textbf{Transparency.} Random CV hides deployment mismatch by mixing periods; a temporal split makes the mismatch visible. In this case study we report both, but we do not propose a mandatory reporting standard: field-level requirements would need multi-site evidence that we do not have.

Our reporting distinguishes four increasingly stringent conditions; the first two are satisfied here, while the latter two are only partially satisfied and are stated as such.

\begin{enumerate}
\item \emph{No future-label leakage.} Features available only after the outcome is known, such as a supersession counter, are inadmissible for event prediction. We satisfy this condition after the controlled diagnostic.
\item \emph{Future-input control.} Every feature is assigned an as-of timestamp. Snapshot covariates such as days since recording and aggregate read counts are legitimate only when frozen at the declared prediction time.
\item \emph{Temporal split.} Training precedes testing in event time. Because event prevalence is bursty, we report test positive rate, event density by period, and censoring rather than a single aggregate score.
\item \emph{Counterfactual audit.} We do not meet this condition: the entity-grounding check is a controlled experiment but underpowered, and the cross-system comparison uses simplified rules. We therefore do not claim transferability.
\end{enumerate}

Our own results satisfy the first rule strictly and the second for the main temporal protocol. The write-time decomposition treats \emph{days since recording} as diagnostic rather than deployable, and the cross-system entity effect is labeled as system-specific. This is the standard we propose: not merely that a paper show a temporal line, but that it disclose where each future-dependent signal enters the pipeline.

The accountability path has three stakeholders. A \emph{memory owner} needs to know why a fact was flagged; the never-read signal can express this as a checkable audit record. A \emph{system operator} needs a reproducible deployment-aligned evaluation before adoption. A \emph{downstream user} needs provenance and an appeal route before deletion. Read telemetry also creates a privacy surface because access counts reveal attention even when content is hidden; deployments should aggregate or pseudonymize counts and disclose retention. If prevalence collapses or telemetry is missing, the baseline should abstain.
\subsection{Limitations}
\label{sec:limitations}

The strongest threat to external validity is the deployment footprint: one production system, one operator, and one governance/logging stack, despite 15 client instances and 22 adapters. We do not claim that the 1.2--1.6 descriptive ratio transfers unchanged. The narrower claim is that event predictions depending on future writes are exposed to temporal leakage unless the protocol proves otherwise.

The cross-system check reinforces this caution. Applying deterministic MemStrata-style rules to our records produced 693 supersessions, compared with 386 natively, with Jaccard overlap 0.120; entity-grounded records were no longer protected at the native level (23.0\% versus 0.1\% supersession). The entity-grounding association is therefore system-specific. A controlled counterfactual was underpowered (estimated power 35\%; about 600 paired cases for 80\% power), so we treat the effect as an implementation hypothesis supported by stratified observations.

The LoCoMo experiments are a bounded negative control, not an event-task replication. Temporal splitting did not change three content-facing tasks: evidence identification (0.638 random versus 0.604 temporal), speaker prediction (0.373 versus 0.380), and temporal reference resolution (0.455 versus 0.439). Because these tasks are not trained to predict future write events over the same record stream, stability is unsurprising. The contrast delimits our claim: temporal validation is necessary where the target depends on the future, but it does not repair post-hoc labels.

Finally, the corpus has a bursty history: the late-September burst produced many procedure-like, entity-free records, and the later documentation phase left long eventless stretches. Temporal estimates are therefore sensitive to the cut-off. Per-agent comparisons are observational; sources differ in workflow, volume, read patterns, and content mix, not merely in NER coverage.

\section{Conclusion}
This case study separates a controlled post-outcome field from ordinary non-stationarity. Removing the former is necessary but not sufficient; deployment-aligned temporal validation is also required. The evaluation protocol can change feature-set rankings. Content-based tasks need the same disclosure discipline, but our LoCoMo checks show they may remain stable under ordinary temporal splitting. For governance, a post-write never-read audit reviews 16.8\% of facts and catches 93.9\% of supersessions at the audited operating point; it is a triage signal, not a predictor, and needs bias and privacy checks. We release the evaluation protocol, annotation codebook, category labels, and analysis scripts. The primary corpus is production-derived; because it contains operational paths, endpoint names, and client-specific configuration traces, we disclose it through the aggregate statistics and de-identified annotations reported here rather than a raw public dump. The annotation labels and scripts allow the categorization and scoring logic to be audited without exposing operational data.

\begin{thebibliography}{20}

\bibitem[Wu et~al.(2024)]{wu2024longmemeval}
Wu, C. et al. LongMemEval: Benchmarking Chat Assistants on Long-Term Interactive Memory. arXiv:2410.10813, 2024.

\bibitem[Chhikara et~al.(2025)]{mem0}
Chhikara, P. et al. Mem0: Building Production-Ready AI Agents with Scalable Long-Term Memory. arXiv:2504.19413, 2025.

\bibitem[Rasmussen et~al.(2025)]{zep}
Rasmussen, R. et al. Zep: A Temporal Knowledge Graph Architecture for Agent Memory. arXiv:2501.13956, 2025.

\bibitem[Yadav(2026)]{memstrata}
Yadav, N. Temporal Validity in Retrieval Memory: Eliminating Stale-Fact Errors for AI Agents over Evolving Knowledge. arXiv:2606.26511, 2026.

\bibitem[Wang et~al.(2026)]{engram}
Wang, L. Less Context, More Accuracy: A Bi-Temporal Memory Engine for LLM Agents. arXiv:2606.09900, 2026.

\bibitem[Anonymous(2026)]{calm}
Anonymous. CALM: Conflict-Aware Lifecycle Memory Governance for Long-Horizon Language Agents. arXiv:2610.38573, 2026.

\bibitem[Tashman(2000)]{tashman2000}
Tashman, L. J. Out-of-sample tests of forecasting accuracy: an analysis and review. \emph{International Journal of Forecasting} 16, 4 (2000), 437--450.

\bibitem[Bergmeir and Benítez(2012)]{bergmeir2012}
Bergmeir, C. and Benítez, J. M. On the use of cross-validation for time series predictor evaluation. \emph{Information Sciences} 191 (2012), 192--213.

\bibitem[Cerqueira et~al.(2020)]{cerqueira2020}
Cerqueira, V., Torgo, L. and Mozetič, I. Evaluating time series forecasting models: an empirical study on performance estimation methods. \emph{Machine Learning} 109 (2020), 1997--2028.

\bibitem[Gama et~al.(2014)]{gama2014}
Gama, J., Žliobaitė, I., Bifet, A., Pechenizkiy, M. and Bouchachia, A. A survey on concept drift adaptation. \emph{ACM Computing Surveys} 46, 4 (2014), Article 44.

\bibitem[Kaufman et~al.(2012)]{kaufman2012}
Kaufman, S., Rosset, S., Perlich, C. and Stitelman, O. Leakage in data mining: formulation, detection, and avoidance. \emph{ACM Transactions on Knowledge Discovery from Data} 6, 4 (2012), Article 15.

\bibitem[Kapoor and Narayanan(2023)]{kapoor2023}
Kapoor, S. and Narayanan, A. Leakage and the reproducibility crisis in machine-learning-based science. \emph{Patterns} 4, 12 (2023), 100804.
\end{thebibliography}

\end{document}











































